\chapter{OBJETIVOS}
\label{cap:objetivos}

% -------------------------------------------------------------
\section{Objetivo Geral}
\label{sec:objetivo-geral}
% -------------------------------------------------------------

Desenvolver um sistema de monitoramento e alerta para barragens de
rejeitos de mineração, baseado em Internet das Coisas (IoT), que
combine a saturação do solo, obtida no local, com a chuva observada e
prevista, e emita alertas em três níveis, de forma local e remota, com
alimentação redundante por bateria, como apoio ao aviso antecipado à
população a jusante.

% -------------------------------------------------------------
\section{Objetivos Específicos}
\label{sec:objetivos-especificos}
% -------------------------------------------------------------

Para alcançar o objetivo geral, foram estabelecidos os objetivos específicos:

\begin{enumerate}[label=\alph*)]

    \item Definir um modelo de alerta com dois limiares, que combine a
    saturação relativa do solo ($S$), obtida pelo higrômetro, e a chuva
    de 72~h ($P_{72}$), formada por 48~h observadas (BNDMET) e 24~h
    previstas (OpenWeatherMap) \cite{mirus2018developing};

    \item Construir o protótipo físico com ESP8266 NodeMCU e a simulação
    no Wokwi com ESP32, com alerta local por LEDs e \textit{buzzer};

    \item Dimensionar por cálculo a bateria de redundância de energia e
    demonstrar o seu funcionamento na simulação;

    \item Desenvolver a plataforma \textit{web}, com página inicial
    (\textit{landing page}), cadastro de moradores, painel do
    administrador e envio de e-mail de alerta (verde, amarelo ou
    vermelho) sobre a necessidade de deslocamento;

    \item Elaborar os fluxogramas e os diagramas de casos de uso e de
    sequência do sistema;

    \item Validar o sistema por cenários de teste que combinem níveis
    de saturação, de chuva e de energia.

\end{enumerate}